% Context from chapters-tex/sec-stochastic-optim.tex; for mathematical reference, not compilation.
conditional expectation, then average over the history. With $e_k=\EE\norm{x_k-x^\star}^2$,
\eql{\label{eq-sgd-proof-2}
e_{k+1}\leq(1-2\mu\tau_k)e_k+\tau_k^2C^2.
}
The claim holds at $k=0$ by the definition of $R$. Set $m=k+2\geq2$. Since $1-2/m\geq0$, the induction hypothesis gives
\eq{e_{k+1}\leq\left(1-\frac2m\right)\frac Rm+\frac R{m^2}
=\frac{R(m-1)}{m^2}\leq\frac R{m+1}.}
The objective bound follows from $L$-smoothness at the minimizer.
\end{proof}

The preceding SGD bound, and the averaging scheme below, give sublinear rates even when $f$ is strongly convex. By contrast, batch gradient descent has a linear rate under the assumptions of Theorem~\ref{thm-gradsec-non-strong-conv}.

Figure~\ref{fig-sgd} compares the objective values $f(x_k)$.
The black dashed curve shows batch gradient descent, with the iteration
axis rescaled to reflect the $n$-fold cost of a batch update.


\begin{figure}
\centering
\begin{tabular}{cc}
\includegraphics[width=.45\linewidth]{ml/sgd/error-sgd-1} &
\includegraphics[width=.45\linewidth]{ml/sgd/error-sgd-2} \\
$f(x_k)$ & $\log_{10}(f(x_k)-f(x^\star))$ 
\end{tabular}
\caption{\label{fig-sgd}
Objective-error evolution for SGD in logistic classification; the dashed curve shows batch gradient descent.
}
\end{figure}




                                                        
\subsection{Stochastic Gradient Descent with Averaging (SGA)}
\label{sec-sga}

Diminishing steps can slow SGD as
$\tau_k$ approaches zero.
 
Averaging reduces fluctuations in the estimate. Run SGD on auxiliary
iterates $( \tilde x_k)_k$,
\eq{
	 \tilde x_{k+1}=\tilde x_k - \tau_k \nabla f_{i(k)}(\tilde x_k)
}
and return their Cesàro average:
\eq{
	x_k \eqdef \frac{1}{k} \sum_{\ell=1}^k \tilde x_{\ell}.
}
This gives stochastic gradient descent with averaging (SGA),
also called averaged SGD.

The average can be maintained without storing the full trajectory,
using the recursive update
\eq{
	x_{k+1}=\frac{1}{k+1}\tilde x_{k+1}+\frac{k}{k+1}x_k.
}


A typical step-size schedule is
\eq{
	\tau_k \eqdef \frac{\tau_0}{1 + \sqrt{k/k_0}}.
}
These steps decrease more slowly, at rate $k^{-1/2}$.


Averaging can reduce sensitivity to the tuning parameters
$(k_0,\tau_0)$, although their choice still affects performance.

                                       

Bach proves that for logistic classification, 
averaging can adapt to local curvature under suitable regularity and step-size assumptions. This is a problem-dependent result, rather than a universal improvement of every constant over SGD.



                                                        
\subsection{Stochastic Averaged Gradient Descent (SAG)}

                                 
For a finite dataset with $n$ samples of total size $n \times p$,
one can trade additional memory for reduced gradient noise
by storing past gradient evaluations. This yields
the stochastic averaged gradient (SAG) method.

Store the latest gradient of each summand in a table $(G^i)_{i=1}^n$, requiring $O(np)$ memory. Maintain its average $g$ as an approximation to the full gradient.

Initialize $G^i=0$ for all $i$ and $g=0$, so that $g=n^{-1}\sum_iG^i$ is the stored average. At each iteration, update the table and then $x$:
\eq{
	x_{k+1} = x_k - \tau g
	\qwhereq
	\choice{
	h\leftarrow\nabla f_{i(k)}(x_k), \\
	g\leftarrow g+\frac1n(h-G^{i(k)}),   \\
	G^{i(k)} \leftarrow h.
	}
}
Because most entries were evaluated at earlier iterates, the table average is generally biased conditional on the current history. Its convergence analysis therefore differs from the unbiased-SGD proof above. SAG uses a fixed step size $\tau$. As for batch gradient descent,
choose $\tau$ on the scale of $1/L$,
where $L$ bounds the Lipschitz constants of the individual gradients $\nabla f_i$; the admissible numerical constant depends on the convergence theorem.

For smooth convex finite sums with a minimizer and a suitable step size, SAG admits an $O(1/k)$ expected objective bound for averaged iterates.
Furthermore, in the presence of strong convexity (for instance after adding a positive quadratic penalty to logistic regression), it has a linear convergence rate,
i.e. 
 \eq{
	\EE( f(x_k) ) - f(x^\star) = O\pa{ \rho^k },
}
for some $0 < \rho < 1$. 

These improvements exploit the finite-sum structure:
SAG stores information about each of the $n$ summands. SGD and averaged SGD
also extend beyond finite $n$ to general
expectations~\eqref{eq-min-int}.

Figure~\ref{fig-compariso-sgd} shows a comparison of SGD, SGA and SAG.

\begin{figure}
\centering
\includegraphics[width=.5\linewidth]{ml/sgd/sg-comparison}
\caption{\label{fig-compariso-sgd}
Evolution of $\log_{10}(f(x_k)-f(x^\star))$ for SGD, SGA and SAG.
}
\end{figure}



